\begin{afterword}
\summaryhead

The post returns to the two figures of ``The Bottom Line.'' The curious inquirer lists all the clues about two boxes and then estimates where the diamond is. The clever arguer writes the conclusion first and then picks clues to support it. The author calls the first procedure rationality and the second rationalization.

The post objects to the word ``rationalization.'' In the author's view it wrongly suggests a kind of rationality, when the process is its opposite. Rationality is described as a forward flow from evidence to conclusion and rationalization as a backward flow from conclusion to selected evidence. The post adds that, for a Bayesian, the expected posterior equals the prior, so an inquirer who already knows the answer has nothing to learn.

The author then criticizes ``Traditional Rationality'' for approving of scientists who set out to prove a pet hypothesis. While granting that such bias drives science forward, the post says it would be better to test hypotheses out of curiosity.

\responsehead

The post has one idea, that reasoning from evidence differs from arguing for a fixed answer, and it states that idea four times. What it adds is a vocabulary: forward flow, backward flow, the Way of Bayes, and ``the Force.'' The additions do not survive inspection.

Its first argument is about a word, and it is wrong. ``Rationalize'' means to make something seem reasonable, and that is what rationalization does. The post calls whoever chose the word a ``fool'' and never checks what the word means or where it came from.

It also leaves out a qualification posted the day before. ``What Evidence Filtered Evidence?'' says that each of the clever arguer's statements ``is valid evidence.'' This post treats the clever arguer as the pure opposite of reason and does not mention it.

Then there is the rival it invents. ``Traditional Rationality'' has no named member and no quoted text; its one line of dialogue is written by the author. The actual tradition, in Popper's formulation, holds that a genuine test is an attempt to refute. The post takes that tradition's central rule, presents it as a correction of the tradition, and gives no credit. What this gives the reader is an identity: members of a new school who see what the old one missed. The price of that identity is a false history of the old school.

The post also leaves the reader where ``The Bottom Line'' did. It says to make sure the flow is not running backwards and gives no way to tell from the inside, where rationalization feels like reasoning.

In short: a post that condemns deciding first and arguing afterwards opens by calling a dictionary word foolish without checking it, and credits to a new school a rule the old one already had.
\end{afterword}
